% Related lecture: SC4061-L12.
% Source: LEC07 Imaging Geometry - Flipping Session, PDF p. 15,
%         Peer-Instruction Question 6, parts 1--3.
% Original formal questions; conversational checks and slide mini-quizzes excluded.
% Source page and calculations verified; no forced page breaks between questions.

\exercisegroup{SC4061-L12-EX}{第 12 讲：相机标定与双目测深}

\examplemeta
  {SC4061-L12-EX}
  {2026-11-03}
  {LEC07 Imaging Geometry -- Flipping Session，文件 p.~15，Question 6}
  {三个分题的题干、公式及答案已核对}
  {学习范围已确认（2026-11-03）}

\exercisequestion{SC4061-L12-E01}{焦距与水平视野角}
\nopagebreak[4]

\textbf{来源：}翻转课堂 p.~15，Question 6.1。
\textbf{对应：}\relatedtopic{SC4061-L12-T02}{DeepFocal}。
\nopagebreak[4]

\textbf{题目。}DeepFocal 使用 $H_\theta=2\arctan(w/(2f))$。
图像宽度为 $w=800\,\mathrm{pixel}$、焦距为 $f=400\,\mathrm{pixel}$，求水平视野角。
当焦距增至 $800\,\mathrm{pixel}$ 而图像宽度不变时，视野角怎样变化？

\begin{checkedsolution}
\begin{align*}
  H_\theta=2\arctan\frac{800}{2\times400}
  &=2\arctan1=\boxed{90^\circ},\\
  H_\theta'=2\arctan\frac{800}{2\times800}
  &=2\arctan0.5\approx\boxed{53.13^\circ}.
\end{align*}
固定图像宽度，增大焦距会缩小视野；视野角与焦距的关系含反正切函数，焦距加倍不会使视野角恰好减半。
\end{checkedsolution}

\exercisequestion{SC4061-L12-E02}{PoseNet 损失权重比较}
\nopagebreak[4]

\textbf{来源：}翻转课堂 p.~15，Question 6.2。
\textbf{对应：}\relatedtopic{SC4061-L12-T02}{位置与朝向的损失权重}。
\nopagebreak[4]

\textbf{题目。}PoseNet 的课程损失为
$L=\|\bm x-\hat{\bm x}\|_2+\beta\|\bm q-\hat{\bm q}\|_2$。
预测 A 的位置误差为 $1.0$、朝向表示误差为 $0.4$；预测 B 的两项误差分别为 $2.0$、$0.1$。
分别在 $\beta=1$ 与 $\beta=10$ 时判断哪个预测的损失更小，并说明权重为何影响选择。

\begin{checkedsolution}
\begin{center}
\begin{tabular}{@{}cccc@{}}
  \toprule
  $\beta$ & $L_A$ & $L_B$ & 较小者 \\
  \midrule
  1 & $1+1(0.4)=1.4$ & $2+1(0.1)=2.1$ & A \\
  10 & $1+10(0.4)=5$ & $2+10(0.1)=3$ & B \\
  \bottomrule
\end{tabular}
\end{center}
A 的位置误差更小，B 的朝向误差更小；提高 $\beta$ 更重视朝向，因此较优预测由 A 变为 B。
题中给出的两项已经是范数，朝向\emph{表示}误差也不能直接读成角度误差。
\end{checkedsolution}

\exercisequestion{SC4061-L12-E03}{单图定位的适用范围}
\nopagebreak[4]

\textbf{来源：}翻转课堂 p.~15，Question 6.3。
\textbf{对应：}\relatedtopic{SC4061-L12-T02}{PoseNet 的场景信息与监督标签}。
\nopagebreak[4]

\textbf{题目。}PoseNet 能从一张图像估计相机位姿，是否说明它不需要场景专属信息，
就能在任意此前未见过的场景中准确定位？解释原因。

\begin{checkedsolution}
不能。单图描述推理时的输入数量；训练仍使用场景图像及对应位姿标签，学习外观与位置、朝向的关系。
SfM 自动生成的标签也属于场景专属信息。陌生场景通常需要相应数据或模型适配。
\end{checkedsolution}
